\subsection{Baza danych}
\label{sec:baza-danych}

Dane aplikacji są przechowywane w bazie PostgreSQL 16 \cite{postgresql},
a dostęp do nich zapewnia Prisma ORM w wersji 7 \cite{prisma}. Schemat bazy
jest zdefiniowany w jednym pliku \path{backend/prisma/schema.prisma}.
Na jego podstawie Prisma generuje klienta z typami TypeScript, więc zapytania
do nieistniejących kolumn i błędy typów wykrywa już kompilator. Prisma 7 nie
korzysta z osobnego silnika zapytań; łączy się z bazą przez sterownik
\texttt{pg} (adapter \texttt{@prisma/adapter-pg}), z limitem 5~sekund na
nawiązanie połączenia.

\subsubsection{Migracje}

Każda zmiana schematu jest zapisywana jako migracja SQL w katalogu
\path{backend/prisma/migrations/} i wersjonowana w repozytorium razem
z kodem. Programista tworzy migrację poleceniem
\verb|npm run db:migrate -w backend -- --name <nazwa>|. W środowisku Docker
migracje stosuje jednorazowa usługa \texttt{migrate}
(\texttt{prisma migrate deploy}), która uruchamia się, gdy baza danych jest
gotowa, a backend startuje dopiero po jej poprawnym zakończeniu. Obraz tej
usługi powstaje z etapu budowania, bo tylko on zawiera narzędzie wiersza
poleceń Prisma. Obraz produkcyjny backendu go nie zawiera.

\subsubsection{Model danych}

Pierwsza migracja (\texttt{init}) tworzy tabelę \texttt{users} z kontami
użytkowników (tabela~\ref{tab:users}). Identyfikatory są typu UUID, więc nie
da się ich odgadnąć, a to ma znaczenie przy adresach dokumentów w API.

\begin{table}[htbp]
\centering
\caption{Tabela \texttt{users}}
\label{tab:users}
\small
\begin{tabular}{|l|l|p{0.45\textwidth}|}
\hline
\textbf{Kolumna} & \textbf{Typ} & \textbf{Opis} \\
\hline
\texttt{id} & \texttt{UUID} & Klucz główny \\
\hline
\texttt{username} & \texttt{VARCHAR(30)} & Login, unikalny \\
\hline
\texttt{email} & \texttt{VARCHAR(254)} & Adres e-mail, unikalny \\
\hline
\texttt{password\_hash} & \texttt{TEXT} & Skrót hasła (argon2), nigdy samo hasło \\
\hline
\texttt{created\_at} & \texttt{TIMESTAMPTZ} & Data utworzenia konta \\
\hline
\texttt{updated\_at} & \texttt{TIMESTAMPTZ} & Data ostatniej zmiany \\
\hline
\end{tabular}
\end{table}

\douzupelnienia{diagram ERD i tabele profilu kierowcy, dokumentów, wersji
i historii zmian; Etapy 5--6}
